\chapter{Podsumowanie}
Pomyślnie zrealizowano wszystkie założone cele projektu. Wytworzono moduły detektora, serwera API oraz serwera webowego.\ Detektor rozpoznaje zdarzenia na nagraniu testowym, prawidłowo anonimizuje twarze wykrytych osób.\ W trakcie prac zidentyfikowano problem związany z rozdzielczością obrazu. Moduł detektora wymagał precyzyjnego dopasowania hiperparametrów do specyfiki kamery testowej, aby zapewnić wysoką skuteczność wykrywania obiektów na rejestrowanych nagraniach.

\section{Dalsze możliwości rozwoju}
Istnieje wiele potencjalnych ścieżek rozwoju, jakie może przyjąć projekt. Możliwe jest m.in.:
\begin{enumerate}
\item Zwiększenie liczby wykrywanych obiektów i zdarzeń;
\item Zoptymalizowanie pętli przetwarzania klatek nagrań;
\item Dodanie mechanizmu uwierzytelniania i autoryzacji w interfejsie użytkownika;
\item Dodanie generowania dodatkowych opisów zdarzeń np.\ za pomocą modeli wizyjno-językowych;
\item Zabezpieczenie punktów końcowych serwera \texttt{API} przed nieupoważnionym dostępem;
\item Zastosowanie obecnego systemu w grze terenowej, w której gracze będą musieli wykonywać określone czynności na kamerach monitoringu. 
\end{enumerate}
Dzięki swojej modularnej budowie system może być rozbudowywany o nowe funkcjonalności, a obecny interfejs graficzny może zostać zastąpiony innym, gdyby zaszła taka potrzeba. Warto byłoby poważnie rozważyć
zabezpieczenie serwera \texttt{API} przed nieupoważnionym dostępem, aby zbierane nagrania z kamer były dodatkowo chronione. Można rozważyć czy gromadzone nagrania nie powinny być przechowywane w datalake'u, aby
zapewnić większą skalowalność i bezpieczeństwo danych. Możliwości rozwoju projektu są szerokie i nie brakuje pomysłów na jego dalszą rozbudowę.